\documentclass[12pt]{article}
\usepackage[T1]{fontenc}
\usepackage[utf8]{inputenc}
\usepackage{lmodern}
\usepackage{textcomp}
\usepackage{enumitem}
\usepackage{geometry}
\geometry{margin=1in}
\begin{document}
\begin{center}
Answer Key - Business Administration - Graduate Level Assessment 15 \
\small (Answers are ordered exactly as in the question paper.)
\end{center}

\section*{Multiple Choice}
\begin{enumerate}[label=\arabic*.]
\item \textbf{Answer:} A
\\ \textit{Options:} A) \$240,000; B) \$40,000; C) \$100,000; D) \$300,000; E) \$280,000
\\ \textit{Note:} The capital of the firm is calculated as total assets minus total liabilities, which in this case is $240,000 (total assets) - $100,000 (total liabilities) = $140,000, but the correct answer provided ($240,000) likely refers to the total asset value, emphasizing the importance of understanding a firm's resources rather than just its net worth.
\item \textbf{Answer:} A
\\ \textit{Options:} A) \$820; B) \$500; C) \$730; D) \$54,600; E) \$2,190
\\ \textit{Note:} The correct answer is $820 because the commission is calculated by multiplying the total sales of $54,600 by the commission rate of 1.5\%, resulting in an income of $54,600 x 0.015 = $819, which rounds to \$820.
\item \textbf{Answer:} A
\\ \textit{Options:} A) \$218.40; B) \$220.00; C) \$234.00; D) \$210.00; E) \$225.60
\\ \textit{Note:} To calculate Mr. Williams' property taxes, multiply the assessed value of $7,800 by the tax rate of $2.80 per $100 (or 0.028), resulting in $7,800 * 0.028 = \$218.40.
\item \textbf{Answer:} A
\\ \textit{Options:} A) \$1,200; B) \$1,215.75; C) \$1,230.50; D) \$1,235.00; E) \$1,242.50
\\ \textit{Note:} The proceeds of $1,200 were calculated by determining the discount on the 90-day note at a rate of 6.5% from the date of discounting, which results in a deduction of $50 from the face value of \$1,250.
\item \textbf{Answer:} B
\\ \textit{Options:} A) 2(1/2)\%; B) 3(1/3)\%; C) 4\%; D) 5\%
\\ \textit{Note:} The effective rate of yield is calculated by dividing the annual dividend of $2 by the stock price of $60, which results in a yield of approximately 3.33\% or 3(1/3)\%.
\end{enumerate}

\section*{True / False}
\begin{enumerate}[label=\arabic*.]
\setcounter{enumi}{5}
\item \textbf{Answer:} False
\\ \textit{Options:} A) True; B) False
\\ \textit{Note:} The statement is false because the selling price of an automobile can vary based on factors such as depreciation, market demand, and pricing strategies, and it is not necessarily fixed at 90\% of its original cost.
\item \textbf{Answer:} True
\\ \textit{Options:} A) True; B) False
\\ \textit{Note:} The statement is true because an annual interest rate of 21.8\% accurately reflects the high-interest terms commonly associated with installment plans for consumer goods like furniture.
\item \textbf{Answer:} True
\\ \textit{Options:} A) True; B) False
\\ \textit{Note:} The statement is true because using the formula for simple interest, \( I = P \times r \times t \), where \( I \) is the interest earned, \( P \) is the principal amount ($1,640), \( r \) is the annual interest rate (4% or 0.04), and \( t \) is the time in years, the calculation shows that $6.56 in interest is generated over 36 days, which is approximately 0.0986 years.
\item \textbf{Answer:} False
\\ \textit{Options:} A) True; B) False
\\ \textit{Note:} PEST analysis includes Political, Economic, Social, and Technological factors, not Environmental or Strategic factors.
\item \textbf{Answer:} False
\\ \textit{Options:} A) True; B) False
\\ \textit{Note:} The statement is false because digital media, with its interactive features and real-time engagement capabilities, often proves to be more effective than print media in showcasing product benefits and energizing advertising messages.
\end{enumerate}

\section*{Long Form Response}
\begin{enumerate}[label=\arabic*.]
\setcounter{enumi}{10}
\item \textbf{Answer:} To calculate the interest paid on a loan of $384.75 taken for 60 days at an interest rate of 6%, we can use the formula for simple interest: Interest = Principal \times Rate \times Time. Here, the principal is $384.75, the annual interest rate is 6\% (or 0.06), and the time period is 60 days, which is equivalent to 60/365 years. Thus, the interest calculation is $384.75 \times 0.06 \times (60/365), resulting in approximately $3.85. Factors influencing the total interest include the principal amount, the interest rate, and the time duration of the loan, which are crucial for financial decision-making as they affect the overall cost of borrowing and cash flow management. Understanding these elements can help individuals and businesses make informed choices regarding loans and investments.
\\ \textit{Note:} The answer correctly applies the simple interest formula by identifying the principal, converting the time period into years, and using the annual interest rate to accurately calculate the interest paid on the loan, while also recognizing the key factors that influence total interest and their relevance to financial decision-making.
\item \textbf{Answer:} Norman Stevens is subject to a property tax calculated based on his home's assessed value and the local tax rate. If his property tax amounts to \$152.31, this figure reflects the application of the assessed value multiplied by the town's tax rate, demonstrating the direct link between property valuation and tax obligations. A portion of this tax is typically allocated to the county government, which plays a crucial role in funding essential services such as public safety, infrastructure, and education. Understanding this allocation is vital for local finance, as it underscores the reliance of municipal budgets on property tax revenues, influencing fiscal policy and community investment decisions. Thus, the property tax levied on Stevens not only impacts his financial responsibilities but also contributes significantly to the broader economic framework of the county.
\\ \textit{Note:} correct because it accurately describes how property tax is calculated based on assessed value and tax rate, and it highlights the allocation of a portion of this tax to the county government, emphasizing its significance in funding vital local services.
\end{enumerate}

\vfill
\noindent\textit{End of Answer Key}
\end{document}